\documentclass[11pt,a4paper]{scrartcl}
\usepackage{iutparakou}
\begin{document}

\fichetitre{Rapport — Jour 4}{Calculabilité : machine de Turing \& machine universelle}
{Binôme : DOSSOU Persévérance \& AGBOGBA Silas \quad$\bullet$\quad Modules : \texttt{formlang/turing.py, utm.py} — App : \texttt{apps/mtu}}

\section*{A. Réponses théoriques}

\subsection*{Q4.1 (0,5) — encodage $\langle M\rangle$}

\paragraph{Schéma.} On numérote (implicitement, par leur nom de chaîne) les états
$Q$ et les symboles $\Gamma$ de $M$. La table $\delta$ est linéarisée comme un
dictionnaire de quintuplets $(q,a)\mapsto(q',b,d)$ avec $d\in\{L,R,S\}$, sérialisé
en JSON (clé textuelle \texttt{"q|a"} $\to$ liste \texttt{[q',b,d]}), auquel on
ajoute l'état initial, l'ensemble des états acceptants, le symbole blanc et
l'ensemble des états de rejet. Le tout est produit par
\texttt{json.dumps(\dots, sort\_keys=True)} : $\langle M\rangle$ est donc une
\textbf{chaîne de caractères unique}.

\paragraph{Injectivité.} Deux machines distinctes (transitions, état initial,
acceptants, blanc ou rejet différents) produisent deux dictionnaires Python
distincts, donc deux chaînes JSON distinctes (JSON est une sérialisation
déterministe ici grâce à \texttt{sort\_keys=True} : pas d'ambiguïté d'ordre) :
\texttt{encode} est injective.

\paragraph{Décodabilité.} \texttt{decode} est la réciproque exacte : elle relit le
JSON, sépare chaque clé \texttt{"q\textbar a"} sur le premier \texttt{"\textbar"}
pour retrouver $(q,a)$, et reconstruit un \texttt{TuringMachine} avec exactement
les mêmes champs. Le \emph{piège} de l'énoncé (« sortie non triée $\Rightarrow$
round-trip non déterministe ») est neutralisé par le tri systématique des clés
avant sérialisation : $\texttt{decode}(\texttt{encode}(M))$ redonne une machine
strictement équivalente à $M$ pour tout mot d'entrée.

\subsection*{Q4.2 (0,5) — la machine universelle}

\paragraph{Description.} $U$ prend en entrée $\langle M\rangle \# w$ (ici, deux
arguments séparés : la description $\langle M\rangle$ et le mot $w$). $U$
\textbf{décode} $\langle M\rangle$ (retrouve $Q,\Gamma,\delta,q_0,F$), puis
\textbf{simule} pas à pas $M$ sur $w$ en utilisant la même boucle générique
\texttt{TuringMachine.run} — $U$ ne réécrit aucune logique propre à $M$, elle
délègue entièrement l'exécution.

\paragraph{Invariant maintenu.} À chaque étape $t$ de la simulation, le ruban de
$U$ (restreint à la zone utile) est exactement le ruban qu'aurait $M$ après $t$
pas d'exécution directe sur $w$, et l'état courant simulé correspond à l'état
réel de $M$ à ce même pas.

\paragraph{Argument de correction (récurrence sur le nombre de pas $t$).}
\emph{Base} ($t=0$) : le ruban initial de la simulation est $w$, comme celui de
$M(w)$ à $t=0$ — vrai par construction. \emph{Hérédité} : si l'invariant est vrai
au pas $t$ (même état, même ruban), alors $U$ lit le même symbole que $M$
lirait, applique la même transition $\delta(q,a)=(q',b,d)$ (car $\delta$ décodée
est exactement celle de $M$), écrit le même symbole, déplace la tête de la même
façon $\Rightarrow$ l'invariant est vrai au pas $t+1$. Par récurrence, $U$
s'arrête en acceptant $w$ ssi $M$ accepte $w$, avec le même ruban final —
exactement ce que vérifie \texttt{test\_utm\_simule\_comme\_la\_machine}.

\subsection*{Q4.3 (0,5) — surcoût (sourcé)}

Simuler une machine multi-ruban par une machine à un seul ruban entraîne un
ralentissement au plus \textbf{quadratique} dans le nombre de pas ; à l'inverse,
simuler $M$ par $U$ (programme enregistré) coûte $O(t\cdot|\langle M\rangle|)$ par
pas simulé, où $t$ est le nombre de pas de $M$ (chaque pas de $M$ nécessite de
reparcourir $\langle M\rangle$ pour retrouver la transition applicable). Hennie
et Stearns (\emph{Two-Tape Simulation of Multitape Turing Machines}, J. ACM
13(4), 1966) montrent qu'un ralentissement seulement \textbf{logarithmique}
(plutôt que quadratique) est possible entre certaines classes de machines
multi-ruban ; aucune borne n'est ici affirmée « de mémoire » sans cette
référence, conformément à la charte d'honnêteté intellectuelle de l'énoncé.

\subsection*{Q4.4 (0,5) — indécidabilité}

\paragraph{Réduction depuis \textsc{halt}.} Le problème \textsc{halt} (« $M$
s'arrête-t-elle sur $w$ ? ») est indécidable (Turing, 1936). Supposons par
l'absurde qu'il existe un décideur $D$ pour « $U$ s'arrête sur
$\langle M\rangle \# w$ ». Comme $U$ simule fidèlement $M$ sur $w$ pas à pas
(Q4.2), $U$ s'arrête sur $\langle M\rangle\#w$ \emph{si et seulement si} $M$
s'arrête sur $w$. On pourrait alors décider \textsc{halt} en construisant
$\langle M\rangle$ (calculable, Q4.1) et en appelant $D$ dessus — contradiction
avec l'indécidabilité de \textsc{halt}. Donc \textbf{« $U$ s'arrête sur
$\langle M\rangle\#w$ » est indécidable}.

\paragraph{Commentaire.} « Universel $\neq$ omniscient » : $U$ peut exécuter
\emph{n'importe quelle} machine décrite, mais elle ne peut pas, en général,
prédire à l'avance si cette exécution va terminer. L'universalité (simuler tout
calcul) et la décidabilité (prédire l'issue de tout calcul) sont deux propriétés
distinctes ; le projet gagne la première au jour 4 mais perd définitivement la
seconde — c'est la leçon annoncée dès la présentation générale de l'énoncé.

\section*{B. Exercices de programmation}

\begin{description}[leftmargin=1.6em,style=nextline]
  \item[E4.1] \texttt{TuringMachine.run} (\texttt{formlang/turing.py}) : ruban
    représenté par un \texttt{dict} \{position $\to$ symbole\} (bi-infini, pas de
    limite de taille), tête = un entier \texttt{pos}. À chaque pas : lecture de
    \texttt{tape.get(pos, blank)}, recherche de $(q,a)$ dans la table ; absence de
    transition $\Rightarrow$ arrêt immédiat (mot rejeté si l'état courant n'est
    pas acceptant, \textbf{sans lever d'exception} — piège de l'énoncé). Un
    compteur \texttt{max\_steps} (option \texttt{trace=True} pour obtenir la trace
    complète pas-à-pas) garantit qu'une boucle infinie lève une erreur explicite
    plutôt que de bloquer indéfiniment. \emph{Test \texttt{test\_trace\_disponible} : passé.}
  \item[E4.2] \texttt{encode}/\texttt{decode} (\texttt{formlang/utm.py}) : voir
    Q4.1 ; \texttt{UniversalTM.run} décode puis délègue à
    \texttt{TuringMachine.run}. \emph{Tests \texttt{test\_round\_trip\_encodage},
    \texttt{test\_utm\_simule\_comme\_la\_machine} : passés.}
  \item[E4.3] Tables \texttt{ADD}/\texttt{SUB} (\texttt{apps/mtu/machines.py}) et
    \texttt{Calculatrice} (\texttt{apps/mtu/calculator.py}) :
    \begin{itemize}
      \item \textbf{ADD} : 3 états. On parcourt $n$ (état \texttt{q0}), le
        \texttt{+} est \textbf{fusionné} en \texttt{1} (état \texttt{q1}), on
        parcourt $m$ jusqu'au blanc, demi-tour, et on \textbf{retire un} \texttt{1}
        (compensation exacte du \texttt{+} converti) avant d'accepter — trace
        conforme à l'exemple de l'énoncé ($2+1=3$ sur \texttt{11+1}).
      \item \textbf{SUB} : appariement par marquage. On cherche un \texttt{1} non
        marqué dans $m$ (marqué \texttt{X}, de gauche à droite, états
        \texttt{q1}/\texttt{q2}), \emph{apparié} à un \texttt{1} non barré de $n$,
        pris le \textbf{plus à droite} (barré \texttt{Y}, états \texttt{q3}/\texttt{q4}) —
        ce choix « droite $\to$ gauche » garantit que les \texttt{1} restants (la
        réponse) forment un bloc \textbf{contigu} à gauche, sans trou. On répète
        jusqu'à épuiser $m$ (nettoyage \texttt{q5}/\texttt{q6} : \texttt{X},
        \texttt{Y}, \texttt{-} effacés, les \texttt{1} non barrés restent) ou
        jusqu'à épuiser $n$ avant $m$ (état \texttt{q\_zero} : $m>n$, tout le
        ruban est effacé, résultat $0$). Les trois points délicats de l'énoncé
        sont traités : gérer ses propres marques (\texttt{q1} saute les
        \texttt{X}, \texttt{q3} saute les \texttt{Y}), barrage droite$\to$gauche
        pour la contiguïté, fin de $m$ détectée par un \textbf{blanc} et non un
        compteur.
      \item \texttt{Calculatrice.addition/soustraction} appellent directement
        \texttt{ADD}/\texttt{SUB} (aucune nouvelle boucle de MT) ;
        \texttt{multiplication} est $m$ additions \emph{composées} ;
        \texttt{division} est une suite de soustractions composées comptant le
        quotient ; \texttt{chainer} applique une liste d'opérations séquentiellement.
    \end{itemize}
    \emph{Tests \texttt{test\_add\_sub\_machines\_directes},
    \texttt{test\_calculatrice\_exhaustif} (49 couples $(n,m)\in\{0,\dots,6\}^2$),
    \texttt{test\_chainage}, \texttt{test\_division\_par\_zero} : passés.}
  \item[E4.4] \texttt{UniversalInterpreter.run(machine, word)} encode
    \texttt{machine} puis délègue à \texttt{UniversalTM} (\texttt{apps/mtu/interpreter.py}) ;
    \texttt{addition\_via\_utm}/\texttt{soustraction\_via\_utm} construisent le
    mot d'entrée unaire puis passent par l'interpréteur universel — aucune
    boucle de ruban n'est réécrite dans \texttt{apps/mtu}, conformément à la
    règle « gate » du jour. \emph{Test \texttt{test\_interpreteur\_universel} : passé.}
\end{description}

\begin{exemple}[Sortie de référence obtenue]
\ttfamily \$ pytest -q tests/test\_turing.py tests/test\_calc.py\\
8 passed
\end{exemple}
\noindent (Sur l'ensemble du dépôt : 25 passed, 3 failed — les 3 échecs restants,
\texttt{test\_nerode\_trois\_classes}, \texttt{test\_pipeline\_word},
\texttt{test\_pipeline\_morpho}, relèvent tous du Jour 5, non encore traité.)

\end{document}
